\documentclass[conference]{IEEEtran}
\makeatletter
\def\ps@IEEEtitlepagestyle{%
  \def\@oddhead{}\def\@evenhead{}}
\makeatother

\usepackage{amsmath,amssymb,amsfonts}
\usepackage{algorithmic}
\usepackage{graphicx}
\usepackage{textcomp}
\usepackage{xcolor}
\usepackage{url}
\usepackage{booktabs}
\usepackage{multirow}
\usepackage{array}
\usepackage{tabularx}
\usepackage{balance}

\begin{document}

\title{MCA Detector: Review-Oriented Discovery of\\
Manipulative Social Coordination}

\author{
  \IEEEauthorblockN{Tsung-Yuan Lin, Jun-Hao Chang, Jia-An Lai}
  \IEEEauthorblockA{
  Department of Computer Science and Information Engineering\\
  Feng Chia University, Taichung, Taiwan\\
  Advisor: Jun-Hong Li}
}

\maketitle

\begin{abstract}
Manipulative activity on social platforms does not always appear as a
single automated account. In cryptocurrency communities, suspicious influence
may emerge from accounts that repeatedly attack the same targets, amplify
similar narratives, or appear in the same discussions within short time
windows. However, highly active legitimate users may also share ideology,
criticize the same opponents, and participate in the same popular threads.
Therefore, treating either an account-level suspiciousness score or shared
negative targets as a final verdict can produce misleading conclusions. This
paper presents MCA Detector, a reproducible and review-oriented pipeline for
discovering suspicious coordination candidates from Reddit cryptocurrency
discussions. The system first uses large language models to analyze post
rhetoric and comment stance, then converts these semantic signals into
account features and multi-layer account graphs. A Manipulative Coordination
Account (MCA) score ranks high-risk seed accounts, while a tiered graph
expansion stage uses co-negative-target relations and supporting graph layers
to form candidate groups. A second-stage temporal verification module then
checks whether candidate group members repeatedly co-appear in the same
thread within short time windows. In the current dataset, MCA Detector
selected 20 seed groups and 172 unique candidate accounts, identifying 13
strong temporal synchrony pairs and 8 moderate temporal synchrony pairs. Case
analysis shows that the harvested group contains both shared negative targets
and robust short-window temporal evidence, while an Odd-Following-247
comparison group contains Stage 1 structure but no reliable Stage 2 timing
evidence. The system is designed to produce review priorities and evidence
tables, not to report supervised detection accuracy or directly accuse
accounts as confirmed bots, cyber troops, or real-world operators.
\end{abstract}

\begin{IEEEkeywords}
social media analysis, coordinated behavior detection, manipulative rhetoric,
large language models, graph expansion, temporal synchrony, Reddit
\end{IEEEkeywords}

\section{Introduction}

Social media platforms have become important environments for public
discussion, financial narratives, and investment sentiment formation. In
cryptocurrency communities, market news, price volatility, policy rumors, and
ideological conflict can quickly produce large volumes of replies. Such
communities are also attractive targets for manipulative activities: a small
number of accounts may attempt to amplify narratives, attack dissenting
voices, or create the appearance of broader agreement.

Detecting this behavior is difficult because manipulation is not equivalent to
automation. Prior social bot detection systems often ask whether an individual
account resembles a bot based on metadata, content, network, and temporal
features \cite{ferrara2016rise,varol2017online}. However, a human-operated
account can still participate in coordinated manipulation, while a highly
active legitimate account may look abnormal without being manipulative. A
second difficulty is that shared targets are ambiguous. If two Bitcoin users
frequently criticize the same anti-Bitcoin author, this may reflect
coordination, but it may also reflect ordinary ideological alignment.

This paper therefore frames the task as suspicious coordination discovery
rather than final account classification. The goal is not to prove that an
account is a bot or a real-world operator. Instead, the goal is to reduce a
large account population into interpretable candidate groups and provide
evidence for human review. The key research question is:

\begin{quote}
How can we identify accounts that deserve further review for manipulative
coordination without overclaiming that all discovered candidates are confirmed
coordinated operators?
\end{quote}

To answer this question, we propose MCA Detector. The system combines
LLM-based semantic analysis, multi-layer account graphs, account-level MCA
seed ranking, graph-based seed expansion, and temporal synchrony verification.
Its central design principle is separation of evidence. MCA score answers
``which accounts are worth starting from?'' Stage 1 graph expansion answers
``who is structurally close to the seed?'' Stage 2 temporal verification
answers ``does this candidate group show short-window co-presence that is more
suspicious than shared ideology alone?''

The contributions of this work are:
\begin{itemize}
  \item We propose a review-oriented pipeline that separates account-level
  suspiciousness, group candidate discovery, and temporal verification.
  \item We convert LLM stance analysis into signed account-to-account
  relations and multiple graph layers, including co-target, co-negative-target,
  trigger-response, and tag-similarity graphs.
  \item We introduce a tiered seed expansion strategy that uses
  co-negative-target overlap as the main discovery signal while requiring
  supporting evidence for noisier layers.
  \item We use same-thread short-window co-appearance as Stage 2 temporal
  evidence and explicitly distinguish temporal strength from temporal
  confidence.
  \item We report case studies and limitations using cautious review language,
  showing how the system ranks suspicious groups without treating every
  candidate as a confirmed cyber troop.
\end{itemize}

\section{Related Work}

\subsection{Social Bot and Account-Level Detection}

Early social bot detection research largely focused on account-level
classification. Ferrara et al. summarized how social bots can affect political,
economic, and public discussions through automated social behavior
\cite{ferrara2016rise}. Varol et al. proposed a framework for estimating
human-bot interactions using metadata, content, sentiment, network, friend,
and temporal features \cite{varol2017online}. These approaches are valuable
for identifying automated or abnormal accounts, but account-level abnormality
does not necessarily imply coordinated manipulation. Our pipeline therefore
uses account-level signals only for seed prioritization and review priority.

\subsection{Group-Level Coordination Detection}

Group-level methods focus on behavioral similarity or synchronized action.
Cresci et al. proposed Social Fingerprinting, which models account behavior
sequences and detects groups of spambots through DNA-inspired representations
\cite{cresci2018social}. Mazza et al. introduced RTbust, which exploits
temporal patterns in retweet time series to detect suspicious botnets
\cite{mazza2019rtbust}. These studies motivate our use of group-level and
temporal evidence rather than relying only on individual account scores.

Coordination network research further represents shared behavioral traces as
networks. Pacheco et al. construct coordination networks from arbitrary shared
behavior traces and demonstrate their use across multiple case studies
\cite{pacheco2021uncovering}. Graham et al. introduce the Coordination Network
Toolkit, a framework for detecting and analyzing coordinated behavior through
weighted, directed multigraphs over multiple behavior traces
\cite{graham2024coordination}. Our work follows this network-oriented view but
adapts it to Reddit reply structure, LLM-derived stance, and candidate-group
review.

\subsection{Stance Analysis and Reddit Troll Detection}

Stance detection identifies whether text is in favor of, against, or neutral
toward a target. The SemEval-2016 stance detection task formalized this
problem for tweets \cite{mohammad2016semeval}. In this project, stance is
used to convert Reddit comments into signed interactions between comment
authors and post authors. This allows us to build negative-target profiles
instead of treating all comments as identical interactions.

Prior Reddit-specific work also motivates cautious interpretation. Saeed et
al. proposed TROLLMAGNIFIER for detecting state-sponsored troll accounts on
Reddit using 335 known Russian-sponsored troll accounts as labeled training
data \cite{saeed2022trollmagnifier}. Their system learns the behavior of known
troll accounts, identifies additional likely trolls, and corroborates the
output using account-level indicators, topic similarity, and temporal
synchronization. Our project differs in that it does not have a labeled troll
set and does not claim supervised troll classification; it produces
review-oriented candidate groups under limited ground truth.

\section{Problem Definition}

Let $A=\{a_1,\ldots,a_n\}$ be the set of accounts, $P$ the set of posts, and
$C$ the set of comments. Each comment $c \in C$ has an author $u(c)$, a target
post $p(c)$, a post author $v(c)$, and a timestamp $t(c)$. The LLM stance
module assigns a feedback label $y(c)$ such as supportive, oppositional,
neutral, mixed, or unclear.

The objective is to produce a ranked list of candidate groups:
\begin{equation}
G_k = (s_k, M_k, E_k, R_k),
\end{equation}
where $s_k$ is a seed account, $M_k$ is a set of candidate members, $E_k$ is
the evidence table for account pairs, and $R_k$ is a review priority label.
Importantly, $R_k$ is not a legal or factual judgment that the members are
confirmed operators. It is a prioritization signal for further review.

We define suspicious manipulative coordination as a pattern in which accounts
show a combination of:
\begin{itemize}
  \item manipulative or attack-oriented content;
  \item shared interaction targets, especially shared negative targets;
  \item sufficient interaction reach or visibility;
  \item abnormal activity patterns;
  \item short-window co-presence in the same discussion context.
\end{itemize}

This definition is intentionally evidence-based. No single signal is treated
as sufficient. In particular, co-negative-target overlap is treated as a
candidate discovery signal, while temporal synchrony is used as a stronger
review signal.

\section{System Architecture}

Fig.~\ref{fig:pipeline} shows the core analysis pipeline. The demo website is
separated from the research pipeline: it reads generated tables and visualizes
the evidence, but it does not compute scores or make final decisions.

\begin{figure}[htbp]
\centering
\fbox{\begin{minipage}{0.94\columnwidth}
\centering
\small
Raw Reddit posts/comments\\
$\downarrow$\\
LLM post rhetoric and comment stance analysis\\
$\downarrow$\\
Account feature matrix and multi-layer adjacency graphs\\
$\downarrow$\\
MCA seed ranking\\
$\downarrow$\\
Stage 1: tiered coordination expansion\\
$\downarrow$\\
Stage 2: temporal synchrony verification\\
$\downarrow$\\
Group summary, candidate table, pair evidence, and roles
\end{minipage}}
\caption{MCA Detector analysis pipeline.}
\label{fig:pipeline}
\end{figure}

\subsection{LLM-Based Semantic Analysis}

The post analysis module produces sentiment and manipulative rhetoric features
such as \texttt{sentiment\_score}, \texttt{manipulative\_rhetoric\_score}, and
\texttt{rhetoric\_tags}. The rhetoric tags include categories such as fear,
urgency, authority claim, bandwagon, overconfidence, emotional amplification,
us-versus-them framing, and call to action.

The comment analysis module determines how a comment responds to the post
author. For a comment written by account $i$ under a post written by account
$j$, the system creates a directed interaction edge:
\begin{equation}
i \rightarrow j.
\end{equation}
Supportive comments contribute positive interactions; oppositional or attacking
comments contribute negative interactions. This design converts unstructured
comment text into account-level graph evidence while preserving interaction
direction and stance.

\subsection{Multi-Layer Account Graphs}

The system constructs several sparse account-level graphs. Table
\ref{tab:graphs} summarizes the major graph layers and their role in the
pipeline.

\begin{table}[htbp]
\centering
\caption{Graph layers and their formal use.}
\label{tab:graphs}
\scriptsize
\setlength{\tabcolsep}{3pt}
\begin{tabularx}{\columnwidth}{p{0.22\columnwidth}X}
\toprule
Graph layer & Role in the system \\
\midrule
Count & Raw comment frequency; baseline and visualization. \\
Signed & Net supportive/oppositional interaction direction. \\
Positive & Supportive interaction context. \\
Negative & Oppositional interaction context. \\
Trigger & Whether one account often responds to another account's posts; support signal. \\
Co-target & Whether two accounts comment on the same targets; support signal. \\
Co-negative & Whether two accounts oppose the same targets; main Stage 1 discovery signal. \\
Tag sim. & Similarity of rhetoric tag profiles; auxiliary support and visualization. \\
\bottomrule
\end{tabularx}
\end{table}

The co-negative-target graph is defined by first building a negative target
profile vector $\mathbf{n}_i$ for each account $i$, where each dimension
corresponds to a target author and the value records the account's negative
interaction strength toward that target. The edge weight is:
\begin{equation}
w_{ij}^{neg} =
\frac{\mathbf{n}_i \cdot \mathbf{n}_j}
{\|\mathbf{n}_i\|\|\mathbf{n}_j\|}.
\end{equation}

This graph is meaningful because manipulative coordination often requires
multiple accounts to pressure or discredit similar targets. However, it is
also limited: ideological communities may organically attack the same targets.
For this reason, $w_{ij}^{neg}$ is used for candidate discovery, not for final
accusation.

\section{MCA Score for Seed Ranking}

MCA score is an account-level review-priority score. It is used to choose
where to start expansion, not to decide whether an account is a confirmed
manipulator. The primary score combines four signal groups:
\begin{equation}
S_{MCA}=0.30S_M+0.35S_C+0.15S_R+0.20S_A,
\end{equation}
where $S_M$ is the manipulative-content signal, $S_C$ is the coordinative
signal, $S_R$ is the interaction-reach signal, and $S_A$ is the automatic
behavior signal.

\begin{table}[htbp]
\centering
\caption{MCA signal groups.}
\label{tab:mca}
\small
\begin{tabular}{p{0.25\columnwidth}p{0.62\columnwidth}}
\toprule
Signal & Features \\
\midrule
Manipulative content & Average manipulative rhetoric score, non-neutral rhetoric ratio, oppositional stance ratio. \\
Coordinative evidence & Co-target strength, co-negative-target strength, trigger-response frequency. \\
Interaction reach & Outgoing volume, incoming attention, interaction breadth. \\
Automatic behavior & Isolation Forest anomaly score transformed into an anomaly percentile. \\
\bottomrule
\end{tabular}
\end{table}

Because large-scale labeled ground truth is unavailable, the weights are
interpretable heuristic weights rather than optimized supervised parameters.
The paper therefore avoids claiming classification accuracy from the MCA score
alone. The score is best understood as a seed ranking function.

The manipulative-content signal mixes post-level and comment-level evidence.
The average rhetoric score and non-neutral rhetoric ratio are computed only
from analyzed posts; accounts with no analyzed posts receive zero for these
two post-level terms. The oppositional stance ratio is computed from analyzed
comments. This choice is conservative because comment-only accounts are not
treated as rhetorically manipulative solely because they attack others, but it
also means that pure comment attackers may receive weaker semantic-content
scores. They can still enter later stages through co-negative-target,
trigger-response, interaction-reach, or temporal evidence.

We performed a lightweight sensitivity check by recomputing MCA rankings under
several weight settings. Table~\ref{tab:sensitivity} shows that moderate
variants preserve many top seeds: the earlier 40/40/10/10 alternative retained
17 of the top 20 accounts and had Spearman correlation 0.983 with the primary
ranking, while a coordination-heavy variant retained 18 of the top 20 accounts
and had correlation 0.986. Removing the manipulative component entirely
reduced stability, supporting the use of MCA as a seed-prioritization score
rather than a final classifier.

\begin{table}[htbp]
\centering
\caption{MCA weight sensitivity. Top-20 is overlap with the primary top-20 list.}
\label{tab:sensitivity}
\scriptsize
\begin{tabular}{p{0.34\columnwidth}rrr}
\toprule
Setting & $\rho$ & Top-20 & harvested \\
\midrule
Primary 30/35/15/20 & 1.000 & 20 & 14 \\
Alt. 40/40/10/10 & 0.983 & 17 & 15 \\
No manipulative & 0.821 & 15 & 11 \\
Low manipulative & 0.911 & 17 & 11 \\
Coordination-heavy & 0.986 & 18 & 11 \\
Automation-heavy & 0.993 & 15 & 19 \\
\bottomrule
\end{tabular}
\end{table}

\section{Stage 1: Tiered Coordination Expansion}

Stage 1 expands each high-MCA seed into a candidate group. The expansion uses
a tiered rule set instead of a single weighted sum so that every included
member has an explicit reason. Table \ref{tab:tier} lists the current rules.

\begin{table}[htbp]
\centering
\caption{Tiered seed expansion rules.}
\label{tab:tier}
\small
\begin{tabular}{p{0.16\columnwidth}p{0.70\columnwidth}}
\toprule
Tier & Inclusion rule \\
\midrule
0 & The seed account itself. \\
1 & Direct co-negative-target edge with weight $\geq 0.20$. \\
2 & Tag similarity $\geq 0.90$ plus structural support from co-negative, co-target, or trigger-response evidence. \\
3 & Trigger-response weight $\geq 0.50$ plus co-negative or tag support. \\
4 & Two-hop co-negative expansion: candidate links to at least two accepted members with co-negative weight $\geq 0.20$. \\
\bottomrule
\end{tabular}
\end{table}

The 2-hop rule is important because a peripheral account may not connect
directly to the seed but may connect to several accepted group members.
However, unrestricted multi-hop traversal can cause candidate explosion.
Therefore, Stage 1 uses only co-negative-target links for 2-hop expansion and
requires at least two links to already accepted members.

The thresholds in this stage are not presented as universal optima. They are
interpretable review settings chosen to keep candidate groups small enough for
manual inspection. For example, the co-negative threshold of 0.20 preserves
shared negative-target structure while avoiding the broad components produced
by raw count or co-target graphs. Full threshold sensitivity across platforms
is left for future work.

The output of Stage 1 includes group members, inclusion reasons, internal
coordination edges, shared negative targets, and account-level feature
summaries. This stage should be interpreted as candidate coordination
discovery. It is not yet a strong claim of shared operation.

\section{Stage 2: Temporal Synchrony Verification}

Stage 2 verifies whether accounts in a candidate group appear in the same post
thread within short time windows. For each account pair $(i,j)$, the system
finds all shared post IDs. For each shared post, it computes all absolute
timestamp differences:
\begin{equation}
\Delta t = |t(c_i)-t(c_j)|.
\end{equation}
The pair-level evidence includes the number of shared posts,
\texttt{within\_5min\_count}, \texttt{within\_30min\_count}, median delay,
minimum delay, and shared thread IDs.

\subsection{Temporal Strength Labels}

The pair-level temporal strength label is:
\begin{itemize}
  \item \textbf{strong temporal sync}: at least one same-thread co-comment
  event occurs within 5 minutes.
  \item \textbf{moderate temporal sync}: at least two same-thread co-comment
  events occur within 30 minutes, or the pair appears in at least three shared
  posts with at least one within-30-minute event.
  \item \textbf{weak temporal overlap}: the pair appears in the same thread
  but lacks short-window synchrony.
  \item \textbf{no temporal sync}: the pair has no same-thread overlap in the
  local comments file.
\end{itemize}

The 5-minute and 30-minute windows are also review-oriented thresholds rather
than statistically optimized cutoffs. The 5-minute window captures near
co-arrival in the same discussion, while the 30-minute window captures repeated
short-window response patterns. These labels are calibrated again by temporal
confidence to avoid treating every short-window event as equally reliable.

\subsection{Temporal Confidence}

Temporal strength alone can overstate a single coincidence. A pair may have
one 5-minute event in a popular thread but no repeated pattern. Therefore, the
system also assigns temporal confidence:
\begin{itemize}
  \item \textbf{robust}: repeated short-window synchrony or multi-post
  short-window evidence with a reasonable median delay.
  \item \textbf{moderate\_review}: useful timing evidence that should be
  reviewed manually.
  \item \textbf{fragile}: evidence is based on a single short-window event or
  has a long typical delay.
  \item \textbf{none}: no usable timing evidence.
\end{itemize}

This separation is important. ``Strong'' describes the presence of a very
short timing event, while ``robust'' describes whether the evidence is stable
enough to trust. A strong pair may still be fragile if it depends on a single
event.

\section{Signal Pruning}

Earlier prototypes included text fingerprint distance and account lifecycle
overlap. Both were removed from formal Stage 2 evidence after testing.

\begin{table}[htbp]
\centering
\caption{Signals tested but removed from formal Stage 2 evidence.}
\label{tab:pruning}
\small
\begin{tabular}{p{0.24\columnwidth}p{0.60\columnwidth}}
\toprule
Signal & Reason for removal \\
\midrule
Text fingerprint distance & In a single-topic Bitcoin community, TF-IDF distance often captures topic similarity instead of shared authorship or shared operation. \\
Lifecycle overlap & Activation windows and account lifetimes did not separate the harvested positive case from independent same-topic users. \\
\bottomrule
\end{tabular}
\end{table}

These fields remain in the output schema for compatibility, but they are
filled as empty values and do not participate in ranking, validation, or demo
interpretation. This pruning decision makes the formal evidence narrower but
more defensible.

\section{Implementation and Outputs}

The pipeline is implemented as separate modules so that each stage can be
re-run independently. The front pipeline reads raw posts and comments, performs
LLM analysis, builds feature matrices, and constructs adjacency graphs. The
back pipeline performs MCA scoring, seed selection, Stage 1 expansion, Stage 2
temporal verification, candidate validation, final group ranking, behavior
profiles, and account role labeling.

The main output tables are:
\begin{itemize}
  \item \texttt{account\_mca\_scores.csv}: account-level MCA scores.
  \item \texttt{tiered\_expansion\_members.csv}: members and inclusion reasons
  for each seed group.
  \item \texttt{stage2\_verification\_evidence.csv}: pair-level temporal
  evidence.
  \item \texttt{candidate\_validation\_table.csv}: account-level review
  priority after merging MCA, expansion, temporal, and behavior metadata.
  \item \texttt{final\_group\_summary.csv}: group-level review priority.
  \item \texttt{account\_role\_table.csv}: simple group-role labels such as
  leader-instigator, comment attacker, comment supporter, and context member.
\end{itemize}

\section{Experimental Results}

The current experiment uses Reddit Bitcoin-related discussion data. The system
scored 49,767 accounts and selected the top 20 MCA accounts as seed accounts.
After Stage 1 expansion, the pipeline produced 20 candidate seed groups, 178
group membership rows, and 172 unique candidate accounts. Stage 2 temporal
verification produced 13 strong temporal synchrony pairs and 8 moderate
temporal synchrony pairs. Among the temporal evidence, 3 pairs reached robust
confidence and 44 pairs were labeled as moderate-review evidence.

\begin{table}[htbp]
\centering
\caption{Overall pipeline output summary.}
\label{tab:overall}
\small
\begin{tabular}{lr}
\toprule
Metric & Value \\
\midrule
Selected seed accounts & 20 \\
Candidate seed groups & 20 \\
Group membership rows & 178 \\
Unique candidate accounts & 172 \\
Strong temporal sync pairs & 13 \\
Moderate temporal sync pairs & 8 \\
Robust temporal pairs & 3 \\
Moderate-review temporal pairs & 44 \\
\bottomrule
\end{tabular}
\end{table}

Table \ref{tab:topgroups} shows the top five groups by final review priority.
The ranking is not an accuracy measurement. It indicates which groups deserve
earlier human inspection. In the table, P1 denotes accounts with the highest
review priority, such as MCA-plus-temporal evidence or MCA-plus-extreme-outlier
evidence; P2 denotes secondary review candidates, such as high-MCA-only or
temporal-only accounts.

\begin{table*}[htbp]
\centering
\caption{Top five candidate groups by review priority.}
\label{tab:topgroups}
\small
\begin{tabular}{lrrrrrrrr}
\toprule
Seed group & Members & P1 & P2 & Tier 1 & Tier 4 & Strong & Moderate & Robust \\
\midrule
lol\_camis & 23 & 2 & 3 & 18 & 4 & 4 & 0 & 0 \\
BtcKing1111 & 11 & 2 & 3 & 6 & 4 & 2 & 1 & 0 \\
tzacPACO & 6 & 2 & 1 & 5 & 0 & 0 & 3 & 1 \\
harvested & 8 & 2 & 0 & 7 & 0 & 1 & 0 & 1 \\
iPurchaseBitcoin & 7 & 2 & 0 & 6 & 0 & 1 & 0 & 1 \\
\bottomrule
\end{tabular}
\end{table*}

\subsection{Lightweight Validation Checks}

Because no large-scale confirmed ground truth is available, we do not report
precision, recall, F1-score, or ROC-AUC. Instead, we add lightweight sanity
checks to test whether each stage contributes different evidence. Table
\ref{tab:stageablation} summarizes the evidence available at each stage. MCA
alone produces only a seed ranking. Stage 1 expands these seeds into 20
candidate groups and 172 unique candidate accounts, but the result may still
reflect ideological alignment. Stage 2 adds 928 pair-level temporal evidence
rows, including 13 strong pairs, 8 moderate pairs, and 3 robust pairs.

\begin{table}[htbp]
\centering
\caption{Ablation-style comparison of pipeline stages.}
\label{tab:stageablation}
\scriptsize
\begin{tabular}{p{0.36\columnwidth}rrrr}
\toprule
Setting & Groups & Accts. & Pairs & S/M/R \\
\midrule
MCA only & 0 & 20 & 0 & 0/0/0 \\
MCA + Stage 1 & 20 & 172 & 0 & 0/0/0 \\
MCA + Stage 1 + Stage 2 & 20 & 172 & 928 & 13/8/3 \\
\bottomrule
\end{tabular}
\end{table}

We also compare candidate-group pairs against a random active-pair baseline.
We sampled 5,000 random account pairs from the local comment file, requiring
each sampled author to have at least five comments and using the same
100-comment-per-author cap as Stage 2. This is a lightweight null baseline,
not a formal statistical significance test. As shown in
Table~\ref{tab:temporalbaseline}, candidate pairs are more likely to share the
same post and more likely to show moderate or robust temporal evidence than
random active pairs.

\begin{table}[htbp]
\centering
\caption{Temporal random baseline.}
\label{tab:temporalbaseline}
\scriptsize
\begin{tabular}{p{0.32\columnwidth}rrrr}
\toprule
Pair set & Same-post & Strong & Moderate & Robust \\
\midrule
Candidate pairs & 45.7\% & 1.40\% & 0.86\% & 0.32\% \\
Random active pairs & 8.5\% & 0.30\% & 0.00\% & 0.02\% \\
\bottomrule
\end{tabular}
\end{table}

\section{Case Studies}

The following cases are selected to illustrate different evidence patterns,
not to claim that they are statistically representative of all candidates.
The harvested group is selected because it has the clearest robust temporal
evidence among the review outputs. Odd-Following-247 is selected because it has
strong Stage 1 shared-target structure but no reliable Stage 2 timing evidence.
BtcKing1111 is included briefly to show how a high-MCA influence-hub candidate
should be interpreted cautiously.

\subsection{The harvested Group}

The harvested group is the clearest high-priority case in the current output.
Stage 1 finds 8 members, including 7 Tier 1 co-negative direct members, 9
internal coordination edges, and 11 shared negative targets. Stage 2 then
finds robust temporal evidence: the pair NectarineDirect936 and harvested
appeared together in 17 shared posts, including 2 events within 5 minutes and
4 events within 30 minutes. This pair is labeled
\texttt{strong\_temporal\_sync} with \texttt{robust} temporal confidence.

This case demonstrates why the two-stage design is useful. Shared negative
targets identify a compact candidate group, while repeated short-window
same-thread co-presence raises review confidence.

\subsection{Odd-Following-247 as a Comparison Group}

Odd-Following-247 is useful as a comparison case because it has Stage 1
structure but no reliable Stage 2 timing evidence. Its group contains 11
members, 7 Tier 1 co-negative direct members, 24 internal coordination edges,
and 15 shared negative targets. However, Stage 2 identifies no strong temporal
sync pairs, no moderate temporal sync pairs, and no robust temporal evidence.
The group has weak temporal overlap, but not short-window synchrony.

This comparison supports the paper's central argument: co-negative-target
structure is useful for candidate discovery, but it should not be treated as
proof of coordinated operation.

\begin{table}[htbp]
\centering
\caption{Case comparison between harvested and Odd-Following-247.}
\label{tab:casecompare}
\scriptsize
\begin{tabular}{p{0.48\columnwidth}rr}
\toprule
Metric & harvested & Odd \\
\midrule
Members & 8 & 11 \\
Tier 1 co-negative members & 7 & 7 \\
Internal coordination edges & 9 & 24 \\
Shared negative targets & 11 & 15 \\
Strong temporal sync pairs & 1 & 0 \\
Moderate temporal sync pairs & 0 & 0 \\
Robust temporal pairs & 1 & 0 \\
Weak temporal overlap pairs & 4 & 40 \\
\bottomrule
\end{tabular}
\end{table}

\subsection{BtcKing1111 as an Influence-Hub Candidate}

BtcKing1111 illustrates a different type of output. It has a high MCA score
of 0.819 and high sub-scores across coordination, interaction reach, and
automatic behavior. Stage 1 expansion produces an 11-member group, and the
role table labels the seed as a leader-instigator because it acts as the group
entry point and has high interaction reach. However, its Stage 2 evidence is
not robust: the group has strong and moderate temporal pairs, but no robust
temporal pairs.

This case should not be described as a confirmed same-operator group. A safer
interpretation is that BtcKing1111 is a high-priority influence-hub candidate:
an account that is visible, structurally connected, and worth manual review,
but not yet proven to control or synchronize other accounts.

\section{Discussion}

\subsection{Why Not Use MCA Alone?}

MCA is useful for ranking suspicious accounts, but it is account-centered. A
high MCA score may indicate a highly active, aggressive, or influential user.
It does not by itself demonstrate that a group of accounts is acting together.
Therefore, MCA is used only to select seeds. The expansion and temporal stages
turn the task from single-account ranking into candidate-group evidence
generation.

\subsection{Why Not Use a Single Graph?}

Each graph layer has a different failure mode. Count and signed graphs can be
dominated by raw activity volume. Co-target graphs are broad because commenting
on the same popular authors is common. Trigger-response graphs can be noisy in
active communities. Tag similarity may be missing for accounts with few posts
and may reflect topic similarity rather than coordination. Co-negative-target
is semantically closer to the suspicious behavior of shared attacking targets,
but it still cannot prove manipulation alone. MCA Detector therefore uses a
multi-layer design: co-negative-target is the main discovery layer, while
other layers provide support, visualization, or review context.

\subsection{What the System Can and Cannot Claim}

The system can claim that it finds accounts and groups with higher review
priority for manipulative coordination risk. It can show why a candidate was
included, which targets were shared, which pairs appeared in the same threads,
and whether the timing evidence is robust or fragile. It cannot claim that
every discovered account is a bot, a cyber troop, or controlled by the same
operator. This distinction is essential for responsible use.

\section{Threats to Validity}

\subsection{Lack of Large-Scale Ground Truth}

The project does not currently have a large labeled dataset of confirmed
coordinated and non-coordinated accounts. Therefore, the results are reported
as review candidates, not precision, recall, or F1-score. The harvested and
comparison cases are used for case-based validation, not for broad statistical
claims.

\subsection{Heuristic Weights and Thresholds}

The MCA weights and Stage 1 thresholds are interpretable heuristics. They were
chosen to create reviewable candidates and avoid candidate explosion, but they
are not optimized supervised parameters. The lightweight weight-sensitivity
check suggests that top seeds are reasonably stable under moderate weight
variants, but full threshold sensitivity for co-negative weights, tag
similarity, trigger response, and temporal windows remains future work.

\subsection{Sparsity of LLM-Derived Content Signals}

LLM-derived manipulative-content features are sparse in the current dataset.
Among 49,767 feature accounts, 23.8\% have at least one analyzed post, 11.7\%
have a nonzero average manipulative rhetoric score, and 10.3\% have non-neutral
rhetoric tags. This does not invalidate MCA for seed selection because
top-ranked seeds generally contain post-level signals, but it limits the use
of manipulative content as a population-wide classifier. Accounts with few or
no analyzed posts may receive weaker semantic-content scores, so downstream
graph and temporal evidence remain necessary.

\subsection{Temporal Coincidence in Popular Threads}

Short-window co-appearance may occur by chance in popular threads. The
temporal confidence layer reduces this risk by separating single-event
evidence from repeated or multi-post evidence. The random active-pair baseline
provides a sanity check, but it does not control for all confounders such as
thread popularity, account activity volume, and overlapping active hours.
Human review remains necessary.

\subsection{Activity Inflation}

High activity can inflate several signals at once: active accounts may have
larger interaction reach, more co-target edges, more co-negative edges, and
more opportunities for same-thread overlap. MCA Detector mitigates this by not
interpreting MCA alone and by separating Stage 1 structure from Stage 2 timing
evidence, but high-activity bias remains a limitation.

\subsection{Platform Specificity}

This experiment focuses on Reddit Bitcoin discussions. Other platforms may not
have the same post-comment structure. Applying the method to PTT, X, Facebook,
or messaging platforms would require redefining action traces, such as
push/boo labels, reposts, hashtags, likes, quote posts, or shared URLs.

\subsection{LLM Reliability}

LLM stance and rhetoric analysis can be affected by prompt design, model
choice, and parsing failures. In this project, LLM outputs are not used as
standalone verdicts. Comment stance labels are relatively auditable because
they describe the relation between a comment and the target post author.
Rhetoric tags are more interpretive and are therefore treated only as
auxiliary semantic features. Each LLM judgment stores a textual rationale,
allowing analysts to inspect why a score or tag was assigned. The extracted
signals are then converted into graph and feature signals and constrained by
structural and temporal evidence. Future work should include prompt
consistency checks and manually audited samples.

\subsection{Demo Scope}

The demo website is an evidence viewer rather than the detector itself. It
reads generated tables from the analysis pipeline and visualizes review
priorities, group relationships, and pair evidence. It does not perform live
classification or compute final account verdicts.

\section{Future Work}

Future work includes four directions. First, we plan to build a small
human-labeled validation set with positive, negative, and noisy cases to
measure review precision more systematically. Second, we plan to expand the
current lightweight sensitivity checks into threshold sensitivity and
permutation-style temporal baselines that better control for thread popularity
and account activity. Third, the pipeline can be extended to additional
platforms by redefining platform-specific action traces; for example, PTT can
use explicit push/boo stance labels. Fourth, the current pipeline can be
packaged as a scheduled monitoring workflow, allowing analysts to periodically
update candidate groups and compare changes over time.

\section{Conclusion}

This paper presented MCA Detector, a review-oriented suspicious social
coordination discovery pipeline. The system combines LLM-based semantic
analysis, multi-layer account graphs, MCA seed ranking, tiered graph expansion,
and temporal synchrony verification. Instead of treating suspiciousness as a
single account-level score, MCA Detector produces candidate groups and
pair-level evidence. The current Reddit Bitcoin experiment selected 20 seed
groups and 172 unique candidate accounts, with 13 strong temporal synchrony
pairs and 8 moderate temporal synchrony pairs. Case studies show that
co-negative-target overlap is useful for discovery but insufficient as a final
claim, while temporal synchrony helps prioritize groups for further review.
The system therefore achieves the practical goal of identifying accounts and
groups that deserve inspection for manipulative coordination risk, while
avoiding overclaims about confirmed bots or real-world operators.

\balance

\begin{thebibliography}{00}

\bibitem{ferrara2016rise}
E. Ferrara, O. Varol, C. Davis, F. Menczer, and A. Flammini, ``The rise of
social bots,'' \emph{Communications of the ACM}, vol. 59, no. 7, pp. 96--104,
2016.

\bibitem{varol2017online}
O. Varol, E. Ferrara, C. A. Davis, F. Menczer, and A. Flammini, ``Online
human-bot interactions: Detection, estimation, and characterization,'' in
\emph{Proc. International AAAI Conference on Web and Social Media}, vol. 11,
no. 1, pp. 280--289, 2017.

\bibitem{cresci2018social}
S. Cresci, R. Di Pietro, M. Petrocchi, A. Spognardi, and M. Tesconi, ``Social
fingerprinting: Detection of spambot groups through DNA-inspired behavioral
modeling,'' \emph{IEEE Transactions on Dependable and Secure Computing},
vol. 15, no. 4, pp. 561--576, 2018.

\bibitem{mazza2019rtbust}
M. Mazza, S. Cresci, M. Avvenuti, W. Quattrociocchi, and M. Tesconi,
``RTbust: Exploiting temporal patterns for botnet detection on Twitter,''
arXiv:1902.04506, 2019.

\bibitem{pacheco2021uncovering}
D. Pacheco, P.-M. Hui, C. Torres-Lugo, B. T. Truong, A. Flammini, and
F. Menczer, ``Uncovering coordinated networks on social media: Methods and
case studies,'' in \emph{Proc. International AAAI Conference on Web and Social
Media}, vol. 15, no. 1, pp. 455--466, 2021.

\bibitem{graham2024coordination}
T. Graham, S. Hames, and E. Alpert, ``The coordination network toolkit: A
framework for detecting and analysing coordinated behaviour on social media,''
\emph{Journal of Computational Social Science}, vol. 7, pp. 1139--1160, 2024.

\bibitem{mohammad2016semeval}
S. Mohammad, S. Kiritchenko, P. Sobhani, X. Zhu, and C. Cherry,
``SemEval-2016 task 6: Detecting stance in tweets,'' in \emph{Proc. 10th
International Workshop on Semantic Evaluation}, pp. 31--41, 2016.

\bibitem{saeed2022trollmagnifier}
M. H. Saeed, S. Ali, J. Blackburn, E. De Cristofaro, S. Zannettou, and
G. Stringhini, ``TROLLMAGNIFIER: Detecting state-sponsored troll accounts on
Reddit,'' in \emph{Proc. IEEE Symposium on Security and Privacy}, pp.
2161--2175, 2022.

\end{thebibliography}

\end{document}
